\documentclass[openany]{ctexbook}
\usepackage{amsmath,amssymb}
\usepackage[hidelinks]{hyperref}
\newcommand{\note}[1]{\footnote{本批译注：#1}}
\newcommand{\Ys}{Y^{\mathrm{sing}}}
\newcommand{\Yss}{Y^{\mathrm{sup\,sing}}}
\title{代数几何基础第五卷预备笔记\\8.12：校订续译}
\author{据用户上传英译编辑稿译出}
\date{2026年11月14日}
\begin{document}
\maketitle
\chapter*{范围及校订原则}
本批连续翻译上传底本 \texttt{EGAVaiello.pdf} 第30页“Remark 8.12.”起至第31页“Proposition 8.13.”标题前的全部可辨文字，包含两段编排备忘及原脚注44--48。下一批从命题8.13标题开始。本批是独立增补，不覆盖此前译稿，尚未并入原对话合订本。

本段仍依赖8.8及8.11的条件。8.7(c)的断裂和8.8完整假设尚未恢复；不把这些备注当作已经补全的证明。可辨原文逐段保留；无法确定的公式、术语和残句标明待核。对已确定的错误才作订正，不以猜测补正文。

本批复核所见原会话仍为第七版918页、第0章至19.6.7；未发现本批范围的交付。网上 Blass 英译本与上传底本逐字节相同，不能当作独立原始见证来恢复缺文。

\chapter*{8.12\quad 备注}
作者从几何上将断言(a)解释如下：对于一个充分一般的超平面截面线性束，束中的每个成员或者光滑，或者其几何奇点都是普通二重点。
\note{此处(a)、(b)指紧接前面的8.11。原文还声称：“事实上，立即可见，因而还可以把(a)说得稍精确些：至多有一个这样的几何奇点。”在目前缺失假设的底本中，不能直接采纳这一加强。上一批所给 $k=\mathbb C$、$X=\mathbb G_m$、$f(t)=[1:t^2:t^4]$ 的非分歧二次覆盖，对 $a\ne0$，相切超平面 $a^2Z_0-2aZ_1+Z_2=0$ 的拉回为 $(t^2-a)^2$，在 $t=\pm\sqrt a$ 两点各产生一个二重零点，故不能仅从现存非分歧假设推出“至多一个”。这里保留作者原断言以供对照，但不声称已验证它。普通二重点与8.8所讨论非退化性的具体对应也需按完整假设核查。}

断言(b)的基本含义是：如果 $\Ys=\Yss$，那么，对于一个充分一般的超平面截面线性束，束中所有成员都是光滑的。\note{此句是对8.11(b)的解释，沿用该推论的特征零背景及相关前提，并非对任意特征的无条件断言。原文下一句“无论是否特征零”修饰的是对这种情形的“例外”评价，不应倒读为已证明任意特征的结论。}
上述等式可在解析意义下表述为：$Y^1$ 上某个可逆模的一个截面 $D$ 消失；底本将该可逆模印作
\[
\omega_{X/k}^{\otimes2}\otimes\mathcal O_{Y'}(1,1).
\]
\note{底本这里同时使用 $Y^1$ 和 $Y'$，且未明确标出 $\omega_{X/k}$ 到相应空间的拉回。此公式按可辨印文保留，记号与定义域待核，不把它视为已恢复的精确线丛公式。}

这第二种情形，无论我们是否处于特征零，都应当完全视为例外情形。〔页顶手迹残段：“簇 $L$ 在 $1\ldots T=L$”；字迹不清。〕\footnote{原脚注44：“字迹不清”。}
用古典语言来说，如果没有弄错，就是说：相对于所考虑的射影浸入，$X$ 是“直纹的”（ruled）。\note{此为作者此处暂用的术语，不能直接等同于通常的直纹簇定义。底本后文8.17又明确表示这里的用法与传统用法不一致；本批不替作者作未经确认的重新定义。}
〔而且，如果我们愿意，〕由8.5及其推论，这里已有把这一术语明确化并加以论证所需的一切；如果你有兴致，也可把它同〔法文：\emph{la taupe}〕联系起来。\note{底本保留法文词组而未解释其指代；本译不猜作某个数学定理或术语。}
例如，若 $\dim X=1$，这意味着 $X$ 是一条直线〔字迹不清〕\footnote{原脚注45：“字迹不清”。}，$x\in X(k)$，所以 $T$ 包含〔缺文〕\footnote{原脚注46：“字迹不清”。}。
\note{“直线”句及随后包含关系在底本中断裂，无法据此确定所需浸入、闭性或其他限制；此处只保存可辨残句，不能将其作为完整已证命题使用。}

(b) 若特征为 $p>0$，通常应当给出这样的例子：$\dim P=2$，$X$ 为非奇异代数曲线，8.6的条件不满足，即 $\Ys=\Yss$，但仍有 $\dim T=r-1$；也就是说，给出 $L/K$ 为有限不可分扩张的例子。我太懒了，不愿构造这些例子，但我不怀疑它们存在。\footnote{原脚注47：“把它做出来。——Blass”。}\note{原文此处确实写“8.6”，并重新出现“(b)”编号；两者均保留，引用关系待核。作者表达的是对例子存在的确信，不是已完成的构造或证明。}

\section*{原稿编排备忘}
在(a)中加一条脚注，指向后面的编号小节或段落；在那里我们证明：如果出现这一例外的“直纹”情形，那么对射影浸入作一个简单的修改，就能重新回到8.8的“一般”情形。\note{这是作者对后文的编排指示，并非本批已经完成该修改或核验该证明的声明。}

本节从8.6到这里的部分，无疑可以独立成一个具有微分性质的小节；而本编号小节的开头，应当与接下来的部分合并成一个讨论例外〔超平面？〕\footnote{原脚注48：“请询问 A.G.”。}之维数的小节。我这里只用了 $\Ys$ 的维数是 $r-1$ 这一事实。

\section*{接续位置}
本批止于上述编排备忘末句。下一处：底本第31页“Proposition 8.13.”标题。未译命题8.13；未恢复8.8缺失条件，未宣称8.9开浸入结论已证明。
\end{document}
